% !TEX program = xelatex
\documentclass[UTF8,10pt]{ctexart}
\usepackage[a4paper,left=1.75cm,right=1.75cm,top=1.7cm,bottom=1.7cm]{geometry}
\usepackage{booktabs,tabularx,array,amsmath,fancyhdr}
\usepackage[table]{xcolor}
\definecolor{navy}{HTML}{17344B}
\definecolor{pale}{HTML}{EAF1F5}
\newcolumntype{Y}{>{\raggedright\arraybackslash}X}
\pagestyle{fancy}\fancyhf{}\fancyfoot[R]{\thepage}\renewcommand{\headrulewidth}{0pt}
\setlength{\parindent}{0pt}\setlength{\parskip}{6pt}
\newcommand{\heading}[1]{\vspace{5pt}{\large\bfseries\color{navy}#1}\par\vspace{3pt}}
\newcommand{\newsection}[1]{\clearpage\heading{#1}}
\newcommand{\smallnote}[1]{{\small\color{black!70}#1}}
\begin{document}

\begin{center}
{\Huge\bfseries\color{navy}驾驶安全事故预测与司机安全评价}\par
\vspace{10pt}{\Large 算法与模型说明}\par
\end{center}
\vspace{12pt}

\heading{赛题目标与总体方案}
依据给定车辆的历史驾驶行为，对固定的500辆车分别完成两项任务：预测未来40天是否发生事故或符合定义的未遂事故；为每名司机生成可解释、可动态更新的安全评分。车辆与司机一一对应，以\texttt{gpsno}作为统一评价标识。模型仅使用组委会提供的车辆画像、风险事件、轨迹和IMU数据。

\begin{center}\small
\begin{tabularx}{\textwidth}{p{3.2cm}Y}
\toprule\rowcolor{pale}任务 & 输出及用途 \\\midrule
事故预测 & \texttt{forecast\_result.csv}逐车给出\texttt{accident}和\texttt{risk\_prob}；后者用于AUC评价和风险排序 \\
司机安全评价 & \texttt{safety\_scores.csv}逐车给出0--100分、排名、等级、六维扣分、原因、可信度及管理建议 \\
\bottomrule
\end{tabularx}
\end{center}

\heading{时间与正式标签}
可用历史数据为2026年6月1日至7月31日；最终预测窗口为8月1日至9月9日，共40天。按照赛题答疑，目标窗口内事件11803至少1条，或事件11804至少2条时，车辆标签为1；其余为0。11803、11804以外的行为事件不构成正式正类。线上任务一AUC仅依据\texttt{risk\_prob}排序计算，\texttt{accident}用于格式核验。

\heading{处理流程}
四类输入先按统一车辆ID、事件时间和日期边界清洗，再形成车辆日级汇总。每个观察窗口仅引用截止日及此前的记录，目标标签仅使用截止日之后的记录。缺失传感器信息保持缺失状态；异常速度、坐标和覆盖率记录质量标识，不将无数据视为安全行为。轨迹暴露量按公里、运行小时和活动天数分别汇总，与各类行为事件的计数口径对应。

\newsection{01 · 任务一：事故预测模型}
\heading{特征设计}
每辆车使用最近14天的行为特征预测随后40天的正式标签。候选特征包括：历史事故及未遂次数、疲劳与分心事件、驾驶稳定性、道路规则、管理配合、轨迹里程与工时、活动天数、事件日分布、近期趋势和传感器质量。计数类特征配合暴露量归一化；对事件稀疏且暴露较低的车辆保留稳健的收缩统计量。

正式模型使用69项稳定特征。其中包含非目标事件的单日集中度、跨行为族同日出现比例，以及其他道路事件的每千公里发生率。对历史阶段覆盖和记录机制变化明显的IMU、路口超速11401及相关聚合特征，不纳入正式预测分支；IMU信息仍用于司机安全评价的佐证和数据可信度。最终预测时只使用截至7月31日已知的信息，不预设每辆车未来40天的行驶量。

\heading{样本构建与模型}
构建8个有完整40天后续标签的历史观察截止日（6月14日至21日），每个截止日对应同一批500辆车。正式目标均按“11803至少1次或11804至少2次”计算。采用浅层LightGBM二分类模型：85棵树、学习率0.035、最大深度2、最多5个叶子、最小叶节点样本45，L2正则15、L1正则0.5。训练权重按同一车辆出现的窗口数做倒数调整，并对正类施加适度权重，减少重叠窗口与类别不平衡的影响。

\heading{风险概率与二分类}
模型先产生风险排序分数，再以车辆分组折外预测的名次为输入拟合正则化Logistic映射，输出\texttt{risk\_prob}\,$\in[0,1]$。该映射保持风险排序方向；在历史样本上得到的概率刻度不等同于已验证的未来真实发生率。\texttt{accident}按历史折外样本选定的阈值生成整数0/1，其取值不参与AUC计算。最终输出严格包含\texttt{gpsno,accident,risk\_prob}三列，每辆车一行。

\newsection{02 · 任务二：评分维度、权重和事件口径}
安全评价覆盖赛题要求的高危历史、司机状态、驾驶平稳性、道路规则、管理配合与IMU佐证。下表权重总和为100\%；事件代码与严重度映射固定在随代码提交的评分规则中。事件码只进入对应维度，避免同一条记录在多个维度重复扣分。

\begin{center}\small
\begin{tabularx}{\textwidth}{p{2.5cm}p{1.35cm}p{1.5cm}Y}
\toprule\rowcolor{pale}维度 & 权重 & 暴露口径 & 事件代码及含义 \\\midrule
历史事故与未遂 & 26.60\% & 每千公里 & 11803事故、11804未遂事故 \\
疲劳与分心 & 13.97\% & 每百小时 & 41001闭眼、41029困倦；41002打哈欠、41003分心、41004电话、41005抽烟、41009低头、41023手机 \\
驾驶稳定性 & 26.60\% & 每千公里 & 30000前碰撞预警、30002/30003车道偏移、30005车距过近、30017压线、60292/60294盲区报警 \\
道路规则风险 & 7.38\% & 每千公里 & 11401路口超速、11402弯道超速、11403危险路段超速、11405匝道超速、11406长下坡超速 \\
管理配合度 & 15.45\% & 每活动天 & 41006摄像头遮挡、41021摄像头角度异常 \\
IMU风险佐证 & 10.00\% & 采样比例 & 加速度/角速度异常及速度差；异常采样点不解释为独立事件次数 \\
\bottomrule
\end{tabularx}
\end{center}

\heading{严重度、时间与暴露量}
每种风险事件设定非负严重度：事故11803为10、未遂11804为7；闭眼疲劳41001与前碰撞预警30000为5，其他行为依风险程度赋予1.5至4.5的固定值。观察窗口为截止日之前14天；事件及其对应的里程、工时或活动天数共同使用14天半衰期，保证分子分母的时间口径一致。IMU按异常采样比例计算辅助证据，不按事件条数累计。全部24类事件代码、权重和暴露单位可在评分规则文件逐项核对。

\newsection{03 · 任务二：安全分计算与可解释性}
对维度$d$，令$n_d$为窗口内严重度与时间衰减后的事件量，$e_d$为同口径暴露量，$\mu_d$、$\tau_d$为固定的收缩先验。先计算风险率和其在固定历史参考分布中的百分位：
\[
r_d=\frac{n_d+\tau_d\mu_d}{e_d+\tau_d},\qquad
P_d(r_d)=\operatorname{Percentile}_{\mathrm{fixed},d}(r_d).
\]
为了使每一项实际扣分对应已观测行为，在相同暴露量下构造零事件背景，再计算事件增量扣分：
\[
B_d=w_dP_d\!\left(\frac{\tau_d\mu_d}{e_d+\tau_d}\right),\quad
C_d=\max\{0,w_dP_d(r_d)-B_d\},\quad
S=100-\sum_{d=1}^{6}C_d.
\]
对无先验收缩的IMU佐证维度，背景项取0。$C_d$均非负，在相同暴露量下随着该维度事件严重度增加而不下降。IMU缺失时该维度不凭空扣分，同时在可信度和保守下界中单独提示数据不足。安全分$S$越高表示观察窗口内可归因的风险行为越少。

\heading{逐项原因与对账}
每辆车列出六维扣分、事件记录数、前三项主要原因及建议动作。事件代码台账按同维度衰减严重度占比分配$C_d$，各代码分配值之和等于维度扣分，各维度扣分之和等于$100-S$。同时保留先验背景与完整参考扣分，便于管理者区分统计估计和实际行为证据。贡献分配是账目归因，不代表某条事件的独立因果效应。

\heading{排名、等级与可信度}
按安全分从高到低排名，车队管理容量依次设为A档20\%、B档40\%、C档30\%、D档10\%；等级表达相对复核优先级。可信度由轨迹覆盖天数、IMU覆盖天数、有效速度/坐标比例、里程、工时与事件数据可用性组成，独立于安全分。低可信度车辆先复核设备或数据，不直接用于自动奖惩。IMU缺失时另列保守分数下界，但不改变公开安全分。

\newsection{04 · 历史验证、动态更新与车队应用}
\heading{历史验证}
任务一采用5折车辆分组验证；同一车辆的所有重叠历史窗口始终留在同一折。验证集使用6月21日截止的500辆车及其后40天正式标签，折均ROC-AUC为0.8445，折内名次合并AP为0.4132。任务二以$100-S$作为风险方向：7月17日截止、随后14天的回顾样本有28个正类，AUC为0.8618、AP为0.2448，前50名高风险车辆覆盖12个正类；6月21日截止、随后40天的回顾样本有55个正类，AUC为0.8385、AP为0.4032。

\begin{center}\small
\begin{tabularx}{\textwidth}{p{3.0cm}p{2.0cm}p{2.0cm}Y}
\toprule\rowcolor{pale}安全等级 & 车辆数 & 正类数 & 7月17日后14天正类比例 \\\midrule
A & 100 & 0 & 0.0\% \\
B & 200 & 2 & 1.0\% \\
C & 150 & 14 & 9.3\% \\
D & 50 & 12 & 24.0\% \\
\bottomrule
\end{tabularx}
\end{center}
\smallnote{以上均为公开历史数据上的回顾或内部验证。40天未来真值尚未提供，绝对概率校准及真实提交成绩需以组委会评测为准；重复历史窗口不能视为新的独立车辆。}

\heading{按天、周、月更新}
每日汇总最近14天的事件与暴露量并重新计算分数；周报展示7天分数变化、等级迁移和待复核事项；月报汇总持续风险维度及培训效果。当前交付含连续8天评分记录及7天分数变化。单独的即时提示显示当日事故/未遂记录或7天内安全分下降至少10分，这些提示不改变安全分及正式预测标签。

\heading{运营动作}
高预测风险且行为评分为C/D者优先在24小时内复核；仅预测高风险者安排7日内核查，仅行为评分偏低者安排针对性培训，常规对象进入30日复盘。低可信度对象先检查数据覆盖和设备。司机可查看自身分数、排名、具体扣分行为和变化趋势；车队管理者可用事件代码台账核对原因。预测概率作为单独的复核依据，不混入安全分。

\heading{可复现交付}
提交包含两项任务的CSV结果、说明PDF以及含数据清洗、特征构造、训练、推理和评分代码的源码包。源码包附有依赖与运行说明、固定评分配置和评估汇总；不包含原始比赛数据。运行环境安装依赖后，可从清洗数据或已生成的标准中间表重新生成两项结果，并按车辆数、列名、取值范围和扣分对账规则完成验收。

\end{document}
